\section{Related Work}

\paragraph{Sparse visual localization.}
Classical structure-based localization matches local image descriptors to a 3D
landmark map and estimates pose with a robust geometric solver. Learned local
features improve repeatability and matching, but the landmark topology is
usually inherited from SfM or selected after reconstruction. \method retains
the sparse deployment interface while moving task-specific learning into map
construction.

\paragraph{Gaussian scene representations.}
3D and 2D Gaussian Splatting optimize explicit primitives for photometric
rendering. Several localization systems render appearance or attach features
to Gaussian primitives. Our formulation instead treats the Gaussian model as a
frozen geometric and visibility scaffold: localization identities and PnP
points are reconstructed from real feature tracks and need not correspond
one-to-one with primitives.

\paragraph{Compact localization maps.}
Map compression commonly ranks or prunes existing 3D points according to
visibility, descriptor quality, or pose coverage. Our topology is reconstructed
jointly with descriptor evidence: track-derived anchors form the core, Gaussian
lineage supplies coverage candidates, and mapping-view self-localization
provides task-specific functional evidence.

\TODO{Replace this concise positioning with citation-grounded discussion before
submission.}
